% ============================================================
%  Course Summary: Quantum Computing and Quantum Machine Learning
%  Graduate Course in Physics
%  Compiled from weekly README content
% ============================================================
\documentclass[aspectratio=169,12pt]{beamer}

% ---------- Theme & colours ---------------------------------
\usetheme{Madrid}
\usecolortheme{dolphin}

\definecolor{qblue}{RGB}{10,60,120}
\definecolor{qcyan}{RGB}{0,150,200}
\definecolor{qgold}{RGB}{210,150,0}

\setbeamercolor{frametitle}{bg=qblue,fg=white}
\setbeamercolor{title}{fg=white}
\setbeamercolor{subtitle}{fg=qcyan}
\setbeamercolor{author}{fg=white}
\setbeamercolor{date}{fg=qgold}
\setbeamercolor{titlelike}{parent=palette primary}
\setbeamercolor{palette primary}{bg=qblue,fg=white}
\setbeamercolor{palette secondary}{bg=qcyan,fg=white}
\setbeamercolor{structure}{fg=qblue}
\setbeamercolor{block title}{bg=qblue,fg=white}
\setbeamercolor{block body}{bg=qblue!8,fg=black}
\setbeamercolor{alerted text}{fg=qgold}

\setbeamertemplate{navigation symbols}{}
\setbeamertemplate{footline}[frame number]

% ---------- Packages ----------------------------------------
\usepackage{amsmath,amssymb}
\usepackage{booktabs}
\usepackage{microtype}
\usepackage{tikz}
\usetikzlibrary{positioning,arrows.meta,shapes.geometric}

% ---------- Macros ------------------------------------------
\newcommand{\ket}[1]{|#1\rangle}
\newcommand{\bra}[1]{\langle#1|}
\newcommand{\week}[1]{\textcolor{qgold}{\textbf{#1}}}
\newcommand{\topic}[1]{\textcolor{qcyan}{\textit{#1}}}

% ============================================================
\title{Quantum Computing and Quantum Machine Learning}
\subtitle{Summary FYS5419/9419}
\author{Morten Hjorth-Jensen}
\institute{Department of Physics, University of Oslo}
\date{Spring semester 2026}
% ============================================================

\begin{document}

% ---- Title slide -------------------------------------------
\begin{frame}[plain]
  \titlepage
\end{frame}

% ---- Course philosophy -------------------------------------
\begin{frame}{Course Philosophy and Structure}
  \begin{columns}[T]
    \column{0.55\textwidth}
      \begin{block}{Two main parts}
        \begin{enumerate}
          \item \textbf{Quantum Computing} (Jan--Mar)\\
                Many-body quantum systems, VQE,\\
                fault-tolerant algorithms
          \item \textbf{Quantum Machine Learning} (Apr--May)\\
                NISQ-era variational methods,\\
                QNNs, QPINNs, Boltzmann machines
        \end{enumerate}
      \end{block}
      \vspace{4pt}
      \begin{alertblock}{Central Algorithm}
        The \textbf{Variational Quantum Eigensolver (VQE)}
        links both halves of the course.
      \end{alertblock}
    \column{0.42\textwidth}
      \begin{block}{Key Themes}
        \begin{itemize}
          \item Hilbert spaces \& quantum states
          \item Entanglement \& measurement
          \item Variational circuits (NISQ)
          \item Fault-tolerant algorithms
          \item Quantum-enhanced ML
          \item Solving ODEs/PDEs with quantum methods
        \end{itemize}
      \end{block}
  \end{columns}
\end{frame}

% ---- Week 1 ------------------------------------------------
\begin{frame}{\week{Week 1 (Jan 19--23)} \ Basic Notions of Quantum Mechanics}
  \begin{columns}[T]
    \column{0.48\textwidth}
      \begin{block}{Mathematical Framework}
        \begin{itemize}
          \item Linear algebra reminder
          \item Hilbert spaces and operators
          \item Bra-ket notation: $\ket{\psi}$, $\bra{\phi}$
          \item Inner products \& norms
        \end{itemize}
      \end{block}
    \column{0.48\textwidth}
      \begin{block}{Quantum States}
        \begin{itemize}
          \item One-qubit systems
          \item Pure states vs.\ mixed states
          \item Composite systems
          \item Operator algebra on Hilbert spaces
        \end{itemize}
      \end{block}
  \end{columns}
  \vspace{8pt}
  \begin{alertblock}{Foundation}
    Establish the mathematical language used throughout
    the entire course.
  \end{alertblock}
\end{frame}

% ---- Weeks 2--3 --------------------------------------------
\begin{frame}{\week{Weeks 2--3 (Jan 26 -- Feb 6)} \ Composite Systems, Density Matrices \& Measurements}
  \begin{columns}[T]
    \column{0.48\textwidth}
      \begin{block}{Week 2: Tensor Products \& Entanglement}
        \begin{itemize}
          \item Spectral decomposition
          \item Measurement formalism
          \item Density matrices $\rho = \sum_i p_i \ket{\psi_i}\bra{\psi_i}$
          \item Entanglement: pure \& mixed states
        \end{itemize}
      \end{block}
    \column{0.48\textwidth}
      \begin{block}{Week 3: Density Matrices in Depth}
        \begin{itemize}
          \item Reduced density matrices
          \item Partial traces
          \item Von Neumann entropy
          \item Entanglement measures
        \end{itemize}
      \end{block}
  \end{columns}
  \vspace{6pt}
  \begin{block}{Key Concept: Entanglement}
    \[
      \ket{\Phi^+} = \frac{1}{\sqrt{2}}\bigl(\ket{00}+\ket{11}\bigr)
      \quad\Rightarrow\quad S(\rho_A) = \log 2
    \]
  \end{block}
\end{frame}

% ---- Week 4 ------------------------------------------------
\begin{frame}{\week{Week 4 (Feb 9--13)} \ Entanglement, Entropies \& Quantum Gates}
  \begin{columns}[T]
    \column{0.48\textwidth}
      \begin{block}{Entanglement \& Entropies}
        \begin{itemize}
          \item Review of density matrices
          \item Entanglement entropies
          \item Separability criteria
        \end{itemize}
      \end{block}
    \column{0.48\textwidth}
      \begin{block}{Quantum Gates (Introduction)}
        \begin{itemize}
          \item Single-qubit gates: $X$, $Y$, $Z$, $H$, $S$, $T$
          \item Two-qubit gates: CNOT, CZ, SWAP
          \item Gate universality
          \item Simple Hamiltonian systems
        \end{itemize}
      \end{block}
  \end{columns}
  \vspace{6pt}
  \begin{alertblock}{Bridge to VQE}
    Understanding gates and Hamiltonians prepares for the VQE algorithm.
  \end{alertblock}
\end{frame}

% ---- Week 5 ------------------------------------------------
\begin{frame}{\week{Week 5 (Feb 16--20)} \ Getting Started with the VQE Algorithm}
  \begin{columns}[T]
    \column{0.55\textwidth}
      \begin{block}{Quantum Gates \& Simple Algorithms}
        \begin{itemize}
          \item Quantum circuit model
          \item Parameterised circuits (ansätze)
          \item Gate decompositions
          \item One- \& two-qubit Hamiltonians
        \end{itemize}
      \end{block}
      \vspace{4pt}
      \begin{block}{VQE Overview}
        \[
          E(\bm{\theta}) = \bra{\psi(\bm{\theta})}H\ket{\psi(\bm{\theta})}
          \;\geq\; E_0
        \]
        Minimise $E(\bm\theta)$ over circuit parameters $\bm\theta$.
      \end{block}
    \column{0.42\textwidth}
      \begin{alertblock}{Project 1}
        Implement VQE for simple\\
        one- and two-qubit Hamiltonians.\\[4pt]
        \textbf{Key tasks:}
        \begin{itemize}
          \item Design ansatz circuit
          \item Evaluate energy expectation
          \item Classical optimisation loop
        \end{itemize}
      \end{alertblock}
  \end{columns}
\end{frame}

% ---- Week 6 ------------------------------------------------
\begin{frame}{\week{Week 6 (Feb 23--27)} \ Implementing VQE: Measurements \& Gradients}
  \begin{columns}[T]
    \column{0.55\textwidth}
      \begin{block}{Advanced VQE Topics}
        \begin{itemize}
          \item \textbf{Parameter-shift rule} for gradients:
                \[
                  \partial_\theta E = \tfrac{E(\theta+\frac\pi2)-E(\theta-\frac\pi2)}{2}
                \]
          \item Adaptive VQE (ADAPT-VQE)
          \item Shot-noise and measurement strategies
          \item Simulation codes and debugging
        \end{itemize}
      \end{block}
    \column{0.42\textwidth}
      \begin{block}{Hamiltonians Explored}
        \begin{itemize}
          \item One-qubit spin models
          \item Two-qubit interacting systems
          \item Mapping physics to qubits
          \item Introduction to the\\
                \textbf{Lipkin model}
        \end{itemize}
      \end{block}
  \end{columns}
\end{frame}

% ---- Weeks 7--8 --------------------------------------------
\begin{frame}{\week{Weeks 7--8 (Mar 2--13)} \ VQE for the Lipkin Model \& Jordan-Wigner}
  \begin{columns}[T]
    \column{0.48\textwidth}
      \begin{block}{Week 7: Lipkin Model VQE}
        \begin{itemize}
          \item Two-qubit Hamiltonian implementation
          \item Full VQE pipeline for the\\
                Lipkin collective model
          \item Benchmarking against exact diagonalisation
          \item Circuit optimisation strategies
        \end{itemize}
      \end{block}
    \column{0.48\textwidth}
      \begin{block}{Week 8: Fermionic Hamiltonians}
        \begin{itemize}
          \item Lipkin model: final results
          \item \textbf{Jordan-Wigner transformation}:
                \[
                  a_j \;\rightarrow\; Z^{\otimes j}\otimes\sigma^-
                \]
          \item Mapping fermionic operators to qubits
          \item Beyond the Lipkin model
        \end{itemize}
      \end{block}
  \end{columns}
  \vspace{4pt}
  \begin{alertblock}{Milestone}
    Project 1 results in a working VQE solver.
  \end{alertblock}
\end{frame}

% ---- Week 9 ------------------------------------------------
\begin{frame}{\week{Week 9 (Mar 16--20)} \ Quantum Fourier Transforms}
  \begin{columns}[T]
    \column{0.55\textwidth}
      \begin{block}{From NISQ to Fault-Tolerant Algorithms}
        \begin{itemize}
          \item Discrete Fourier Transform review
          \item \textbf{Quantum Fourier Transform (QFT)}:
                \[
                  \mathrm{QFT}\ket{j} = \frac{1}{\sqrt{N}}
                  \sum_{k=0}^{N-1} e^{2\pi i jk/N}\ket{k}
                \]
          \item Circuit decomposition: $O(n^2)$ gates vs.\ $O(n\,2^n)$
        \end{itemize}
      \end{block}
    \column{0.42\textwidth}
      \begin{block}{QFT Circuit Structure}
        \begin{itemize}
          \item Hadamard gates
          \item Controlled phase rotations $R_k$
          \item Bit-reversal permutation
          \item Exponential speedup over classical FFT
        \end{itemize}
      \end{block}
  \end{columns}
\end{frame}

% ---- Weeks 10--11 ------------------------------------------
\begin{frame}{\week{Weeks 10--11 (Mar 23 -- Apr 10)} \ QFT, Quantum Phase Estimation \& HHL}
  \begin{columns}[T]
    \column{0.48\textwidth}
      \begin{block}{Quantum Phase Estimation (QPE)}
        \begin{itemize}
          \item QFT as subroutine
          \item Estimating eigenvalues of $U$:
                \[
                  U\ket{u} = e^{2\pi i\varphi}\ket{u}
                \]
          \item Circuit construction \& precision
          \item Applications to chemistry \& physics
        \end{itemize}
      \end{block}
    \column{0.48\textwidth}
      \begin{block}{HHL Algorithm (Linear Systems)}
        \begin{itemize}
          \item Solving $A\bm{x}=\bm{b}$ on a quantum computer
          \item Uses QPE as subroutine
          \item Exponential speedup (sparse, well-conditioned $A$)
          \item Implementation and codes
          \item Limitations and caveats
        \end{itemize}
      \end{block}
  \end{columns}
\end{frame}

% ---- Week 12 -----------------------------------------------
\begin{frame}{\week{Week 12 (Apr 13--17)} \ HHL Algorithm: Codes \& Implementation}
  \begin{columns}[T]
    \column{0.55\textwidth}
      \begin{block}{Deep Dive into HHL}
        \begin{itemize}
          \item Full circuit implementation
          \item Ancilla qubit rotation: eigenvalue inversion
          \item Amplitude encoding of $\bm{b}$
          \item Uncomputation and measurement
          \item Error analysis and success probability
        \end{itemize}
      \end{block}
    \column{0.42\textwidth}
      \begin{alertblock}{HHL Complexity}
        \[
          \mathcal{O}\!\left(\log N \cdot \kappa^2 \cdot \frac{1}{\epsilon}\right)
        \]
        vs.\ classical $\mathcal{O}(N^3)$\\[4pt]
        where $\kappa$ = condition number,\\
        $\epsilon$ = precision,\\
        $N$ = matrix dimension.
      \end{alertblock}
    \end{columns}
    \vspace{4pt}
    \begin{block}{Transition to Quantum Machine Learning}
      With QFT, QPE and HHL covered, we now turn to the
      \textbf{NISQ-era QML} methods for the remainder of the course.
    \end{block}
\end{frame}

% ---- Week 13 -----------------------------------------------
\begin{frame}{\week{Week 13 (Apr 20--24)} \ QAOA: Quantum Approximate Optimisation}
  \begin{columns}[T]
    \column{0.55\textwidth}
      \begin{block}{QAOA Algorithm}
        \begin{itemize}
          \item Combinatorial optimisation problems
          \item Alternating cost and mixer unitaries:
                \[
                  \ket{\bm\gamma,\bm\beta}
                  = U_M(\beta_p)\,U_C(\gamma_p)\cdots
                    U_M(\beta_1)\,U_C(\gamma_1)\ket{+}^{\otimes n}
                \]
          \item Classical outer-loop optimisation
          \item Depth-$p$ approximation ratio
        \end{itemize}
      \end{block}
    \column{0.42\textwidth}
      \begin{block}{Implementation Focus}
        \begin{itemize}
          \item Mapping problems to Ising Hamiltonians
          \item Circuit construction
          \item Gradient evaluation
          \item Benchmarking on toy instances
        \end{itemize}
      \end{block}
  \end{columns}
  \vspace{4pt}
  \begin{alertblock}{Connection}
    QAOA is the optimisation counterpart to VQE and serves as
    a gateway into Quantum Machine Learning methods.
  \end{alertblock}
\end{frame}

% ---- Week 14 -----------------------------------------------
\begin{frame}{\week{Week 14 (Apr 27 -- May 1)} \ Quantum Machine Learning: Foundations}
  \begin{columns}[T]
    \column{0.55\textwidth}
      \begin{block}{QML Basics}
        \begin{itemize}
          \item What is QML? Taxonomy of approaches
          \item Data encoding strategies:\\
                amplitude, angle, IQP, etc.
          \item Parameterised quantum circuits as models
          \item Expressibility and entangling capacity
          \item Barren plateaus and trainability
        \end{itemize}
      \end{block}
    \column{0.42\textwidth}
      \begin{block}{Quantum Neural Networks}
        \begin{itemize}
          \item QNN architecture overview
          \item Layered variational circuits
          \item Loss functions \& gradients
          \item Links to classical deep learning
          \item Connections to QAOA
        \end{itemize}
      \end{block}
  \end{columns}
\end{frame}

% ---- Weeks 15--16 ------------------------------------------
\begin{frame}{\week{Weeks 15--16 (May 4--15)} \ QPINNs: Solving Differential Equations}
  \begin{columns}[T]
    \column{0.48\textwidth}
      \begin{block}{Physics-Informed QNNs}
        \begin{itemize}
          \item Classical PINNs review
          \item \textbf{Quantum PINNs (QPINNs)}:
                replace neural network with\\
                a parameterised quantum circuit
          \item Encoding PDEs as loss functions
          \item Automatic differentiation on circuits
        \end{itemize}
      \end{block}
    \column{0.48\textwidth}
      \begin{block}{Applications \& Implementation}
        \begin{itemize}
          \item Solving ODEs and PDEs:\\
                Schrödinger, heat, Poisson
          \item Boundary condition encoding
          \item Gradient computation via\\
                parameter-shift rule
          \item Discussion of codes
        \end{itemize}
      \end{block}
  \end{columns}
  \vspace{4pt}
  \begin{alertblock}{Key Idea}
    QPINNs marry the physics-informed training paradigm with
    the expressive power of quantum circuits to solve differential
    equations on near-term hardware.
  \end{alertblock}
\end{frame}

% ---- Week 17 -----------------------------------------------
\begin{frame}{\week{Week 17 (May 18--22)} \ Quantum Boltzmann Machines}
  \begin{columns}[T]
    \column{0.55\textwidth}
      \begin{block}{Quantum Boltzmann Machines (QBM)}
        \begin{itemize}
          \item Classical RBM / BM recap
          \item Quantum generalisation:\\
                thermal states as model distributions
          \item Training via quantum log-likelihood gradients
          \item Transverse-field Hamiltonians as model parameters
          \item Restricted QBM (RQBM) architecture
        \end{itemize}
      \end{block}
    \column{0.42\textwidth}
      \begin{block}{Connections \& Outlook}
        \begin{itemize}
          \item Links to VQE and variational circuits
          \item Generative quantum models
          \item Applications in physics \& optimisation
          \item Open research problems
        \end{itemize}
      \end{block}
  \end{columns}
  \vspace{4pt}
  \begin{alertblock}{Course Conclusion}
    QBMs illustrate how quantum statistical mechanics and
    machine learning merge
  \end{alertblock}
\end{frame}

% ---- Course map: Part 1 ------------------------------------
\begin{frame}{Course Road Map --- Part 1: Quantum Computing}
  \begin{center}
  \small
  \renewcommand{\arraystretch}{1.28}
  \begin{tabular}{@{}llp{8.0cm}@{}}
    \toprule
    \textbf{Weeks} & \textbf{Dates} & \textbf{Topics} \\
    \midrule
    1     & Jan 19--23     & Basic quantum mechanics, Hilbert spaces, qubits \\
    2--3  & Jan 26--Feb 6  & Composite systems, density matrices, entanglement \\
    4     & Feb 9--13      & Entanglement entropies, quantum gates \\
    5--6  & Feb 16--27     & VQE: algorithm, measurements, gradients (Project 1) \\
    7--8  & Mar 2--13      & VQE for Lipkin model, Jordan-Wigner transform \\
    9     & Mar 16--20     & Quantum Fourier Transform \\
    10--11& Mar 23--Apr 10 & QPE and HHL algorithm \\
    12    & Apr 13--17     & HHL codes and implementation \\
    \bottomrule
  \end{tabular}
  \end{center}
\end{frame}

% ---- Course map: Part 2 ------------------------------------
\begin{frame}{Course Road Map --- Part 2: Quantum Machine Learning}
  \begin{center}
  \small
  \renewcommand{\arraystretch}{1.28}
  \begin{tabular}{@{}llp{8.0cm}@{}}
    \toprule
    \textbf{Weeks} & \textbf{Dates} & \textbf{Topics} \\
    \midrule
    13    & Apr 20--24     & QAOA: quantum approximate optimisation \\
    14    & Apr 27--May 1  & QML foundations, Quantum Neural Networks \\
    15--16& May 4--15      & QPINNs: solving differential equations \\
    17    & May 18--22     & Quantum Boltzmann Machines and course summary \\
    \bottomrule
  \end{tabular}
  \end{center}
  \vspace{6pt}
  \begin{alertblock}{Overarching Goal}
    Equip participants with the theoretical foundations and practical
    implementation skills to design, analyse, and run quantum algorithms
    on near-term and fault-tolerant hardware.
  \end{alertblock}
\end{frame}

% ---- Closing -----------------------------------------------
\begin{frame}{Summary: What we have done}
  \begin{columns}[T]
    \column{0.48\textwidth}
      \begin{block}{Part 1 -- Quantum Computing}
        \begin{itemize}
          \item Formalism of quantum mechanics
          \item Entanglement \& density matrices
          \item Variational algorithms (VQE)
          \item Fault-tolerant algorithms:\\
                QFT, QPE, HHL
          \item Physical models: Lipkin, Pairing
        \end{itemize}
      \end{block}
    \column{0.48\textwidth}
      \begin{block}{Part 2 -- Quantum ML}
        \begin{itemize}
          \item QAOA for optimisation
          \item Quantum Neural Networks
          \item Physics-Informed QNNs\\
                (QPINNs)
          \item Quantum Boltzmann Machines
          \item Connections to classical ML
        \end{itemize}
      \end{block}
  \end{columns}
  \vspace{8pt}
  \begin{alertblock}{Overarching Goal}
    Equip course participants with the theoretical foundations
    and practical implementation skills to design, analyse, and run
    quantum algorithms on near-term and fault-tolerant hardware.
  \end{alertblock}
\end{frame}

% ============================================================
%  PERSPECTIVES — QUANTUM COMPUTING: FUTURE DIRECTIONS
% ============================================================

% ---- Perspectives: section divider -------------------------
\begin{frame}[plain]
  \begin{center}
    \vspace{1.4cm}
    {\Huge\bfseries\color{qblue} Perspectives}\\[0.5em]
    {\Large\color{black} Future Directions in Quantum Computing}\\[0.3em]
    {\normalsize\color{gray} Algorithms, Hardware, and Scientific Applications}
  \end{center}
\end{frame}

% ---- Perspectives 1: NISQ & fault-tolerant algorithms ------
\begin{frame}{Perspectives --- Algorithms: NISQ and Fault-Tolerant}
  \begin{columns}[T]
    \column{0.48\textwidth}
      \begin{block}{Near-Term (NISQ) Algorithms}
        \begin{itemize}\small\itemsep1pt
          \item \textbf{Improved VQE ans\"{a}tze:}
                hardware-efficient, UCC and ADAPT-VQE,
                symmetry-preserving circuits
          \item \textbf{Quantum error mitigation:}
                zero-noise extrapolation,
                probabilistic error cancellation
          \item \textbf{Barren plateau remedies:}
                layerwise training, local cost
                functions, better initialisation
          \item \textbf{Quantum advantage benchmarks:}
                tasks where NISQ devices
                outperform classical methods
        \end{itemize}
      \end{block}
    \column{0.48\textwidth}
      \begin{block}{Fault-Tolerant Algorithms}
        \begin{itemize}\small\itemsep1pt
          \item \textbf{Quantum simulation:}
                Trotterisation, qubitisation,
                linear combination of unitaries (LCU)
          \item \textbf{Improved QPE:}
                randomised and Bayesian phase
                estimation with fewer qubits
          \item \textbf{Beyond HHL:}
                quantum linear algebra (QSVM,
                quantum PCA), dequantisation
          \item \textbf{Quantum walks and search:}
                Grover speedups, walk-based
                graph algorithms, Monte Carlo
        \end{itemize}
      \end{block}
  \end{columns}
\end{frame}

% ---- Perspectives 2: Physical sciences applications --------
\begin{frame}{Perspectives --- Applications in the Physical Sciences}
  \begin{columns}[T]
    \column{0.48\textwidth}
      \begin{block}{Quantum Chemistry and Nuclear Physics}
        \begin{itemize}\small\itemsep1pt
          \item \textbf{Electronic structure:}
                VQE and QPE for ground and excited
                states beyond coupled-cluster reach
          \item \textbf{Nuclear structure:}
                ab initio calculations using Lipkin,
                pairing, and shell-model Hamiltonians
          \item \textbf{Quantum advantage:}
                fault-tolerant simulation of strongly
                correlated systems (e.g.\ FeMoco)
          \item \textbf{Imaginary-time evolution:}
                thermal states and finite-temperature
                properties on quantum hardware
        \end{itemize}
      \end{block}
    \column{0.48\textwidth}
      \begin{block}{Condensed Matter and High-Energy Physics}
        \begin{itemize}\small\itemsep1pt
          \item \textbf{Lattice gauge theories:}
                quantum simulation of QED, QCD,
                and the Hubbard model
          \item \textbf{Topological phases:}
                preparation and detection of
                topological order and anyons
          \item \textbf{Quantum gravity:}
                holographic duality (AdS/CFT)
                encoded on quantum circuits
          \item \textbf{Quantum sensing:}
                Heisenberg-limited parameter
                estimation, magnetometry,
                gravimetry
        \end{itemize}
      \end{block}
  \end{columns}
\end{frame}

% ---- Perspectives 3: QML, hardware, open challenges --------
\begin{frame}{Perspectives --- Quantum ML, Hardware, and Open Challenges}
  \begin{columns}[T]
    \column{0.54\textwidth}
      \begin{block}{Quantum Machine Learning Frontiers}
        \begin{itemize}\small\itemsep1pt
          \item \textbf{Quantum generative models:}
                quantum GANs, quantum diffusion,
                and QBMs beyond the RQBM
          \item \textbf{Quantum reinforcement learning:}
                RL on quantum hardware for
                quantum control and
                agentic scientific workflows
          \item \textbf{Classical--quantum hybrid:}
                tensor networks and classical
                circuit simulation as a
                guide for algorithm design
          \item \textbf{Certification:}
                classical shadow tomography
                and cross-entropy benchmarking
        \end{itemize}
      \end{block}
    \column{0.43\textwidth}
      \begin{block}{Hardware and Open Challenges}
        \begin{itemize}\small\itemsep1pt
          \item \textbf{Qubit platforms:}
                superconducting, trapped-ion,
                photonic, neutral-atom arrays
          \item \textbf{Error correction:}
                surface and color codes;
                logical fault-tolerant qubits
          \item \textbf{Scalability:}
                how many logical qubits for
                quantum advantage on
                real-world problems?
        \end{itemize}
      \end{block}
      \begin{alertblock}{Key Insight}
        Quantum advantage requires both
        better algorithms \emph{and}
        better hardware in tandem.
      \end{alertblock}
  \end{columns}
\end{frame}

\end{document}
